\documentclass[11pt, a4paper]{article}
\usepackage[T1]{fontenc}
\usepackage[utf8]{inputenc}
\usepackage[a4paper, top=2.5cm, bottom=2.5cm, left=2.5cm, right=2.5cm]{geometry}
\usepackage{booktabs}
\usepackage{tabularx}
\usepackage{enumitem}
\usepackage[hidelinks]{hyperref}
\usepackage{parskip}
\usepackage{graphicx}
\usepackage{amsmath}
\usepackage{float}
\setlength{\parindent}{0pt}
\pagestyle{plain}

\begin{document}

% ------------------------------------------------------------------------------
% PORTADA
% ------------------------------------------------------------------------------
\begin{center}
{\Large \textbf{Pontificia Universidad Javeriana}}\\[0.3cm]
Facultad de Ciencias Economicas y Administrativas\\[0.3cm]
Estadistica -- 2026-II\\[1cm]
{\LARGE \textbf{Publica vs privada: ?`que tipo de universidad agrega mas valor?}}\\[0.5cm]
{\large Informe Final del Proyecto}\\[1cm]
\textbf{Grupo 5}\\[0.3cm]
Reyes Rodriguez Jose Miguel\\
Santamaria Caballero Isabella\\
Bermudez Caicedo Gabriela\\[1cm]
27 de noviembre de 2026
\end{center}

\newpage
\tableofcontents
\newpage

% ------------------------------------------------------------------------------
\section{Introduccion}
% ------------------------------------------------------------------------------
% Contexto: debate sobre calidad de educacion publica vs privada en Colombia.
% Pregunta: las IES oficiales agregan mas o menos valor que las no oficiales,
% controlando por el nivel de entrada del estudiante?
% Concepto de valor agregado (delta) como medida de lo que aporta la universidad.

% ------------------------------------------------------------------------------
\section{Datos y Metodologia}
% ------------------------------------------------------------------------------
% Fuente: cruce Saber 11 -- Saber Pro (ICFES).
% Archivos: cruce\_saber11\_saberpro.parquet, saber11.parquet, saber\_pro.parquet.
% Variable dependiente: delta = scale(PUNT\_GLOBAL\_sbpro) - scale(PUNT\_GLOBAL\_sb11).
% Predictores: IES\_tipo (Oficial vs No Oficial, simplificado de INST\_ORIGEN\_sbpro),
%   PUNT\_GLOBAL\_sb11, FAMI\_ESTRATOVIVIENDA\_sb11, INST\_CARACTER\_ACADEMICO\_sbpro.
% Modelo OLS unico con diagnosticos completos.

% ------------------------------------------------------------------------------
\section{Analisis Exploratorio (Fase 1)}
% ------------------------------------------------------------------------------
% Distribucion de delta por tipo de IES.
% Histograma y boxplot del valor agregado.
% Estadisticas descriptivas por tipo de IES y caracter academico.

% ------------------------------------------------------------------------------
\section{Intervalos de Confianza y Pruebas de Hipotesis (Fase 2)}
% ------------------------------------------------------------------------------
% IC para la media de delta por tipo de IES.
% Prueba t para diferencia de medias delta oficial vs no oficial.
% Prueba de Wilcoxon como alternativa no parametrica.

% ------------------------------------------------------------------------------
\section{Modelo de Regresion (Fase 3)}
% ------------------------------------------------------------------------------
% Modelo: delta $\sim$ IES\_tipo + PUNT\_GLOBAL\_sb11 + FAMI\_ESTRATOVIVIENDA\_sb11
%   + INST\_CARACTER\_ACADEMICO\_sbpro.
% Tabla de coeficientes con IC al 95\%.
% Predicciones: delta predicho para IES oficial vs no oficial al P50 de SB11.

% ------------------------------------------------------------------------------
\section{Diagnosticos y Robustez}
% ------------------------------------------------------------------------------
% Graficos de residuales (residuales vs ajustados, QQ-plot, Scale-Location).
% Test de Breusch-Pagan (homocedasticidad).
% VIF (multicolinealidad).
% Distancia de Cook (observaciones influyentes).
% Errores robustos HC1: comparacion con errores estandar clasicos.

% ------------------------------------------------------------------------------
\section{Resultados y Discusion}
% ------------------------------------------------------------------------------
% Magnitud y significancia del coeficiente de IES\_tipo.
% Comparacion de predicciones publica vs privada.
% Reflexion sobre calidad educativa y valor agregado institucional.

% ------------------------------------------------------------------------------
\section{Conclusiones}
% ------------------------------------------------------------------------------
% Respuesta a la pregunta de investigacion.
% Limitaciones del estudio.
% Implicaciones para la discusion sobre financiamiento de educacion publica.

% ------------------------------------------------------------------------------
\section*{Referencias}
% ------------------------------------------------------------------------------
% ICFES, literatura sobre calidad educativa publica vs privada en Colombia.

% ------------------------------------------------------------------------------
\section*{Anexo: Codigo R}
% ------------------------------------------------------------------------------
% Incluir codigo R utilizado (fase3\_grupo5.R) o referenciarlo.

\end{document}
